\documentclass[a4paper,12pt]{article}

\usepackage[T2A]{fontenc}
\usepackage[utf8]{inputenc}
\usepackage[russian]{babel}
\usepackage{amsmath,amssymb,mathtools}
\usepackage{helvet}
\renewcommand{\familydefault}{\sfdefault}
\usepackage{icomma}
\usepackage[a4paper,margin=20mm]{geometry}
\usepackage{microtype}
\usepackage{tikz}
\usetikzlibrary{arrows.meta,positioning,shapes.geometric}
\usepackage{tabularx,booktabs}
\usepackage{enumitem}
\usepackage{parskip}
\usepackage{fancyhdr}
\usepackage{setspace}
\usepackage{xcolor}

\definecolor{accent}{HTML}{8B3A62}
\definecolor{darkaccent}{HTML}{5B2440}
\definecolor{textgray}{HTML}{2F2F2F}
\definecolor{softgray}{HTML}{6D6D6D}

\geometry{headheight=15pt}
\onehalfspacing
\setlength{\parindent}{0pt}
\setlist[itemize]{leftmargin=1.5em,itemsep=0.25em,topsep=0.35em}
\setlist[enumerate]{leftmargin=1.7em,itemsep=0.55em,topsep=0.4em}

\pagestyle{fancy}
\fancyhf{}
\lhead{\small\color{softgray}Алгебраические дроби}
\rhead{\small\color{softgray}Екатерина Скелли}
\cfoot{\small\color{softgray}\thepage}
\renewcommand{\headrulewidth}{0.3pt}

\newcommand{\topic}[1]{%
  \vspace{0.45em}
  {\large\bfseries\color{darkaccent}$\square$\;#1}\par
  \vspace{0.25em}
}

\newcommand{\keyidea}[1]{%
  \begin{center}
  \begin{minipage}{0.92\linewidth}
  \centering\bfseries\color{darkaccent}#1
  \end{minipage}
  \end{center}
}

\tikzset{
  startstop/.style={ellipse,draw=darkaccent,line width=0.8pt,minimum width=34mm,minimum height=11mm,align=center,font=\small},
  process/.style={rectangle,rounded corners=2mm,draw=accent,line width=0.8pt,text width=115mm,minimum height=13mm,align=center,font=\small,inner sep=4mm},
  arrow/.style={-{Stealth[length=3mm]},line width=0.9pt,draw=darkaccent}
}

\begin{document}

\begin{titlepage}
\thispagestyle{empty}
\centering
\vspace*{2.2cm}

{\Large\color{softgray}Математика. 8 класс\par}
\vspace{1.2cm}

{\fontsize{30}{36}\selectfont\bfseries\color{darkaccent}Чек-лист\par}
\vspace{0.45cm}
{\LARGE\bfseries Алгебраические дроби\par}
\vspace{0.25cm}
{\Large Разложение на множители и сокращение\par}

\vspace{2.0cm}
{\large Екатерина Скелли\par}
\vspace{0.5cm}
{\normalsize \today\par}

\vfill
{\small\color{textgray}Лёвин Артём Александрович эксклюзивно для Екатерины Скелли\par}
\vspace{0.7cm}

\begin{tikzpicture}[x=1cm,y=1cm]
\draw[accent,line width=1.2pt]
  (0,0) .. controls (3,0.7) and (7,-0.7) .. (10,0)
        .. controls (12.8,0.7) and (15.6,-0.4) .. (16.6,0.15);
\end{tikzpicture}
\end{titlepage}

\section*{Главная идея занятия}
\keyidea{Чтобы сократить алгебраическую дробь, сначала превращайте суммы и разности в произведения. Сокращать можно только общие множители.}

\topic{1. Вынесение общего множителя}
Если все слагаемые имеют общий множитель, вынесите его за скобки:
\[
 x^2-5x=x(x-5),
 \qquad
 -3-5x=-(3+5x).
\]
Минус тоже удобно выносить как общий множитель. Это особенно важно, когда нужно получить одинаковые множители в числителе и знаменателе.

\topic{2. Разность квадратов}
\[
 u^2-v^2=(u-v)(u+v).
\]
Примеры:
\[
4x^2-9y^2=(2x-3y)(2x+3y),
\]
\[
25-9x^2=(5-3x)(5+3x).
\]
Сумма квадратов $u^2+v^2$ по этой формуле не раскладывается.

\topic{3. Квадрат суммы и квадрат разности}
\[
(u+v)^2=u^2+2uv+v^2,
\qquad
(u-v)^2=u^2-2uv+v^2.
\]
Полезно уметь двигаться в обе стороны:
\[
9x^2+12xy+4y^2=(3x+2y)^2,
\]
\[
x^2-6x+9=(x-3)^2.
\]
Проверка среднего слагаемого: оно должно быть равно $\pm 2uv$.

\topic{4. Знак перед дробью и «переворот» разности}
\[
-\frac{A}{B}=\frac{-A}{B}=\frac{A}{-B}.
\]
Если поменять порядок слагаемых в разности, появляется минус:
\[
2y-3x=-(3x-2y).
\]
Это позволяет согласовать знаки у множителей перед сокращением.

\clearpage
\section*{Сокращение алгебраических дробей}

\topic{5. Сокращаем множители, а не слагаемые}
Если числитель и знаменатель разложены на множители, одинаковые множители можно сократить:
\[
\frac{(x-3)(x+1)}{x(x-3)}=\frac{x+1}{x},
\qquad x\ne0,\;x\ne3.
\]
Запись $x-3$ исчезает после сокращения, однако ограничение $x\ne3$ остаётся: оно пришло из исходного знаменателя.

\topic{6. ОДЗ записываем по исходному выражению}
Перед сокращением отметьте все значения, при которых исходные знаменатели обращаются в ноль. После преобразований эти запреты сохраняются.

Например,
\[
\frac{-3x-2y}{9x^2-4y^2}
=\frac{-(3x+2y)}{(3x-2y)(3x+2y)}
=-\frac{1}{3x-2y},
\]
где
\[
3x-2y\ne0,\qquad 3x+2y\ne0.
\]

Ещё один типичный случай:
\[
\frac{-6x-4y}{9x^2+12xy+4y^2}
=\frac{-2(3x+2y)}{(3x+2y)^2}
=-\frac{2}{3x+2y},
\]
при $3x+2y\ne0$.

\topic{7. Разложение группировкой}
Если четыре слагаемых можно разбить на пары с одинаковой скобкой, используйте группировку:
\begin{align*}
 a^2+3ab-2a-6b
 &=a(a+3b)-2(a+3b)\\
 &=(a+3b)(a-2).
\end{align*}

\clearpage
\topic{8. Полный пример из занятия}
Упростим
\[
\frac{a^2+3ab-2a-6b}{a^2-9b^2}\cdot
\frac{a-3b}{a^2-4}.
\]
Сначала ОДЗ:
\[
a\ne3b,\qquad a\ne-3b,\qquad a\ne2,\qquad a\ne-2.
\]
Разложим всё на множители:
\begin{align*}
&\frac{(a+3b)(a-2)}{(a-3b)(a+3b)}\cdot
\frac{a-3b}{(a-2)(a+2)}\\[2mm]
&\hspace{25mm}=\frac{1}{a+2}.
\end{align*}
Ответ имеет более простой вид, однако все ограничения исходного выражения сохраняются.

\section*{Мини-памятка: частые ошибки}
\begin{itemize}
  \item Не сокращайте части суммы: сначала требуется разложение на множители.
  \item Не забывайте средний член в $(u\pm v)^2$: это $\pm2uv$.
  \item Не путайте $u^2-v^2$ и $u^2+v^2$.
  \item При перестановке членов разности меняйте знак: $u-v=-(v-u)$.
  \item ОДЗ берётся из исходного выражения и не исчезает после сокращения.
  \item Перед финальным ответом проверьте: всё ли, что Вы сократили, действительно было множителем целиком.
\end{itemize}

\section*{Самопроверка перед завершением}
\begin{itemize}
  \item[$\square$] Я вижу общий множитель и умею его вынести.
  \item[$\square$] Я узнаю разность квадратов.
  \item[$\square$] Я узнаю полный квадрат трёхчлена.
  \item[$\square$] Я умею разложить выражение группировкой.
  \item[$\square$] Я записываю ОДЗ до сокращения.
  \item[$\square$] Я сокращаю только множители.
  \item[$\square$] Я умею согласовать знаки у противоположных разностей.
\end{itemize}

\clearpage
\section*{Алгоритм: как упростить алгебраическую дробь}
\thispagestyle{plain}
\vspace{0.4cm}
\begin{center}
\begin{tikzpicture}[node distance=7mm]
\node[startstop] (start) {Начало};
\node[process,below=of start] (odz) {Запишите ОДЗ по всем исходным знаменателям};
\node[process,below=of odz] (factor) {Разложите числители и знаменатели на множители: общий множитель, разность квадратов, квадрат суммы/разности, группировка};
\node[process,below=of factor] (sign) {Если множители отличаются только порядком разности, вынесите минус: $u-v=-(v-u)$};
\node[process,below=of sign] (cancel) {Сократите только одинаковые множители};
\node[process,below=of cancel] (finishcalc) {Перемножьте оставшиеся множители и приведите запись к простому виду};
\node[process,below=of finishcalc] (check) {Проверьте, что ограничения ОДЗ не потерялись после сокращения};
\node[startstop,below=of check] (stop) {Готово};

\draw[arrow] (start) -- (odz);
\draw[arrow] (odz) -- (factor);
\draw[arrow] (factor) -- (sign);
\draw[arrow] (sign) -- (cancel);
\draw[arrow] (cancel) -- (finishcalc);
\draw[arrow] (finishcalc) -- (check);
\draw[arrow] (check) -- (stop);
\end{tikzpicture}
\end{center}
\clearpage

\clearpage
\section*{Тренировочный блок --- Путь А}
\textbf{10 заданий.} Выполняйте разложение на множители явно. В заданиях с дробями указывайте ограничения исходного выражения.

\begin{enumerate}
  \item Разложите на множители:
  \[
  x^2-6x.
  \]

  \item Разложите на множители:
  \[
  16a^2-25b^2.
  \]

  \item Представьте в виде квадрата двучлена:
  \[
  9x^2+24xy+16y^2.
  \]

  \item Упростите:
  \[
  \frac{-4x-8y}{4x^2+16xy+16y^2}.
  \]

  \item Упростите:
  \[
  \frac{x^2-16}{x^2+8x+16}.
  \]

  \item Упростите:
  \[
  \frac{(p-q)^2}{(p+q)^2}\cdot
  \frac{5(p+q)^2}{2(p-q)^2}.
  \]

  \item Упростите:
  \[
  \frac{x^2-9}{x+2y}\cdot
  \frac{4x+8y}{x^2+3x}.
  \]

  \item Упростите, при необходимости согласовав знаки разностей:
  \[
  \frac{-6x+8y}{9x^2-16y^2}.
  \]

  \item Упростите:
  \[
  \frac{a^2+4ab-3a-12b}{a^2-16b^2}\cdot
  \frac{a-4b}{a^2-9}.
  \]

  \item Упростите:
  \[
  \frac{m^2-1}{m+2n}\cdot
  \frac{6m+12n}{m^2+m}.
  \]
\end{enumerate}

\clearpage
\section*{Ответы}
\renewcommand{\arraystretch}{1.35}
\begin{center}
\begin{tabularx}{\linewidth}{@{}c>{\raggedright\arraybackslash}X>{\raggedright\arraybackslash}X@{}}
\toprule
\textbf{№} & \textbf{Ответ} & \textbf{Ограничения для дробей}\\
\midrule
1 & $x(x-6)$ & ---\\
2 & $(4a-5b)(4a+5b)$ & ---\\
3 & $(3x+4y)^2$ & ---\\
4 & $-\dfrac{1}{x+2y}$ & $x+2y\ne0$\\
5 & $\dfrac{x-4}{x+4}$ & $x\ne-4$\\
6 & $\dfrac52$ & $p+q\ne0$, $p-q\ne0$\\
7 & $\dfrac{4(x-3)}{x}$ & $x+2y\ne0$, $x\ne0$, $x\ne-3$\\
8 & $-\dfrac{2}{3x+4y}$ & $3x-4y\ne0$, $3x+4y\ne0$\\
9 & $\dfrac{1}{a+3}$ & $a\ne4b$, $a\ne-4b$, $a\ne3$, $a\ne-3$\\
10 & $\dfrac{6(m-1)}{m}$ & $m+2n\ne0$, $m\ne0$, $m\ne-1$\\
\bottomrule
\end{tabularx}
\end{center}

\vfill
\begin{center}
{\small\color{textgray}Лёвин Артём Александрович эксклюзивно для Екатерины Скелли}
\end{center}

\end{document}
